\section{What the monitor sees on development weeks}
\label{sec:results}

This section summarises the development weeks; Appendix~\ref{sec:results-full} gives the full results and Section~\ref{sec:holdout} re-tests the final pipeline on untouched weeks. Colluders are $K=50$ volume-matched accounts sending about 1{,}640 queries a week. Weeks, not windows, are the independent unit: zero false alarms in 960 windows from 24 weeks bounds the rate at about 12\%, not 0.3\%. Run-to-run variation (0.79--0.93 for one cell across four runs) exceeds single-run intervals, so we compare detectors within a run.

\paragraph{Calibration breaks first.} Thresholds set once on early weeks fired on 12\% of later null weeks against a nominal 1\%, and building the null by resampling individual accounts fails for the same reason: it erases benign coordination. Resampling \emph{whole} past weeks against a lagged reference works. Over whole weekly populations only the two-scale test stayed calibrated under rolling calibration (0.7\% false alarms), against 2.9--14.2\% for every other detector (Appendix Figure~\ref{fig:calib-full}).

\begin{table}[t]
\caption{Weekly paired detection at deployed thresholds (nominal 1\%), whole weekly populations, $K=50$ volume-matched colluders, 24 development weeks. The last row is the realised false-alarm rate of each detector.}
\label{tab:deploy}
\begin{center}
\footnotesize\setlength{\tabcolsep}{4pt}
\begin{tabular}{lccccc}
\toprule
Campaign & Two-scale & Global stream & MMD & Spectral & Per-account / HC / Fraudar \\
\midrule
MetaMathQA (generated math) & \textbf{0.93} & 0.71 & 0.60 & 0.19 & $\le$0.01 \\
Magicoder (generated code) & \textbf{0.62} & 0.46 & 0.32 & 0.00 & 0.00 \\
Alpaca (generated, broad) & 0.01 & 0.00 & 0.01 & 0.00 & 0.00 \\
Narrow cluster of real prompts & 0.12 & 0.04 & \textbf{0.27} & 0.00 & 0.00 \\
Random real prompts & 0.01 & 0.00 & 0.01 & 0.00 & 0.00 \\
\midrule
Realised false alarms & 0.7\% & 4.1\% & 11.6\% & 7.8\% & 7.5 / 5.3 / 2.9\% \\
\bottomrule
\end{tabular}
\end{center}
\end{table}

\paragraph{Coordinated, dataset-like campaigns are catchable.} Topically focused generated datasets were caught in most weeks, while no per-account monitor caught anything (Table~\ref{tab:deploy}). Broad datasets and real user prompts were not---exactly what the barrier predicts. When the two-scale test fired, the flagged group was almost entirely colluders (precision 0.94--1.00) but held only 12--31\% of them: a pointer for review, not a full list. Its edge is calibration rather than ranking: at oracle-matched 1\% false alarms it caught 0.80 of MetaMathQA campaigns against 0.70 for the global stream and 0.58 for MMD, but the global stream ranks slightly better overall (AUC 0.97 against 0.88).

\paragraph{It detects topic shift, not extraction.} We designed the two-scale test on the idea that it would detect deduplicated generated data. Two ablations rejected that idea. Coordinated \emph{human-written} GSM8K questions were caught more often (0.94) than generated MetaMathQA (0.79), while simpler source-shift statistics on the same group---a dataset-versus-chat classifier, topic concentration---caught nothing at 1\%. And benign groups sweeping the GSM8K test set from 3 or 10 accounts triggered the test in 79--80\% of weeks, as often as distillation.

\paragraph{Padding, and a fix that looks at directions.} Mixing natural prompts into every colluder account cut two-scale detection from 0.85 to 0.21--0.24. Four repairs that pool topic cells across accounts ranked padded campaigns better but never reached a usable false-alarm rate. What worked was asking \emph{why} padding works: it shrinks the \emph{size} of each colluder's deviation but not its \emph{direction}. With conformal thresholds, hollowed selection caught padded campaigns 0.46--0.48 of the time, against 0.00--0.02 for selection on $C$, with no false alarm in 960 windows. Once topicality and drift are removed, benign accounts in these weeks were barely more coordinated than chance: the top eigenvalue of $R_c$ was 1.14 times that of shuffled weeks.

\paragraph{Scale: selection fails before the barrier does.} At colluder shares of 1--2.5\% no detector was both calibrated and powerful. In populations pooled from account-weeks across the year (Appendix Figure~\ref{fig:scale-full}), two-scale power fell from 1.00 to 0.00 as the colluder share went from 20\% to 0.5\%, yet the \emph{true} colluder set still scored above the threshold at every size: the detector stopped finding the right accounts. Multi-outlier selection closed that gap with oracle thresholds (93--100\% of unpadded campaigns up to $10^4$ accounts), but its best deployed calibration reached 5\% false alarms, not 1\%, and pooling probably understates same-week benign coordination. \textbf{For deployment at realistic colluder shares, our result is negative.} Query monitoring is much stronger as a complement to metadata linkage: given 30\% of a campaign's colluders, expanding along $R_c$ recovered 96--100\% of the rest (oracle thresholds, known campaign size).

\paragraph{Attacks on the monitor.} Ordering prompts greedily against the two-scale statistic cut its detection from 0.80 to 0.31 without any padding. Giving the most distinctive queries to five sacrificial accounts left selection on $C$ flagging those five and none of the other 45 colluders. And ramping a campaign up through the eight calibration weeks poisoned the reference and the thresholds, cutting detection from 0.95 to 0.12---the boiling-frog attack on anomaly detectors \citep{rubinstein2009antidote}.
